%%%%%%%%%%%%%%%%%%%%%%%%%%%%%%%%%%%%%%%%%%%%%%%%%%%%%%%%%%%%%%%%%%%%%%%%%%%
\chapter{llama-server and llama-embedding}
\label{ch:llama-server}
%%%%%%%%%%%%%%%%%%%%%%%%%%%%%%%%%%%%%%%%%%%%%%%%%%%%%%%%%%%%%%%%%%%%%%%%%%%

%%%%%%%%%%%%%%%%%%%%%%%%%%%%%%%%%%%%%%%%%%%%%%%%%%%%%%%%%%%%%%%%%%%%%%%%%%%
\section{What llama-embedding Does}
%%%%%%%%%%%%%%%%%%%%%%%%%%%%%%%%%%%%%%%%%%%%%%%%%%%%%%%%%%%%%%%%%%%%%%%%%%%

\texttt{llama-embedding} is the llama.cpp command-line tool for generating embeddings
directly from a model instead of generating text. Conceptually it:

\begin{enumerate}
    \item Loads an embedding-capable model
    \item Tokenizes the input text
    \item Runs a forward pass through the model
    \item Extracts a vector representation from the model output
    \item Optionally normalizes or formats that vector for downstream use
\end{enumerate}

Instead of returning next-token probabilities or generated text, it returns a numeric vector
that represents the semantic meaning of the input.

%%%%%%%%%%%%%%%%%%%%%%%%%%%%%%%%%%%%%%%%%%%%%%%%%%%%%%%%%%%%%%%%%%%%%%%%%%%
\section{What Embedding Generation Means at the Model Level}
%%%%%%%%%%%%%%%%%%%%%%%%%%%%%%%%%%%%%%%%%%%%%%%%%%%%%%%%%%%%%%%%%%%%%%%%%%%

A transformer processes input tokens into hidden states. An embedding model then turns those
hidden states into a single fixed-size vector. That vector is what you store in a vector index
and compare during retrieval.

The key challenge is that token sequences have variable length but vector search wants a
fixed-size output. So the tool needs a pooling strategy.

%%%%%%%%%%%%%%%%%%%%%%%%%%%%%%%%%%%%%%%%%%%%%%%%%%%%%%%%%%%%%%%%%%%%%%%%%%%
\section{Why Pooling Matters}
%%%%%%%%%%%%%%%%%%%%%%%%%%%%%%%%%%%%%%%%%%%%%%%%%%%%%%%%%%%%%%%%%%%%%%%%%%%

Pooling is the method used to collapse token-level hidden states into one embedding vector.
Common pooling strategies:

\begin{description}
    \item[\texttt{last}] The final token representation (used by Qwen3 embedding models)
    \item[\texttt{mean}] Average of all token representations
    \item[\texttt{cls}] The special \texttt{[CLS]} token representation (BERT-style)
\end{description}

For Qwen3 embedding GGUF models, the official card explicitly requires:

\begin{Verbatim}
--pooling last
\end{Verbatim}

This matters because using the wrong pooling mode can degrade retrieval quality even if the
model ``runs.'' For Qwen3 embedding models, \texttt{last} is not just an implementation
detail---it is part of using the model correctly.

%%%%%%%%%%%%%%%%%%%%%%%%%%%%%%%%%%%%%%%%%%%%%%%%%%%%%%%%%%%%%%%%%%%%%%%%%%%
\section{llama-embedding vs llama-server \texttt{-{}-embedding}}
%%%%%%%%%%%%%%%%%%%%%%%%%%%%%%%%%%%%%%%%%%%%%%%%%%%%%%%%%%%%%%%%%%%%%%%%%%%

These serve different roles.

\textbf{\texttt{llama-embedding}:}

\begin{itemize}
    \item One-shot CLI tool
    \item Useful for local experiments
    \item Useful for benchmarking
    \item Useful for generating vectors offline in scripts or batch jobs
\end{itemize}

\textbf{\texttt{llama-server -{}-embedding}:}

\begin{itemize}
    \item Long-running service
    \item Exposes an HTTP endpoint (\texttt{/v1/embeddings})
    \item Better for OpenAI-compatible clients and apps like Open WebUI
    \item Better when embeddings need to be requested repeatedly
\end{itemize}

The rough split is: use \texttt{llama-embedding} to inspect and test; use
\texttt{llama-server -{}-embedding} to serve a real application.

%%%%%%%%%%%%%%%%%%%%%%%%%%%%%%%%%%%%%%%%%%%%%%%%%%%%%%%%%%%%%%%%%%%%%%%%%%%
\section{Why Embedding Models Are Often Served Separately}
%%%%%%%%%%%%%%%%%%%%%%%%%%%%%%%%%%%%%%%%%%%%%%%%%%%%%%%%%%%%%%%%%%%%%%%%%%%

A chat model server and an embedding model server usually have different jobs:

\begin{itemize}
    \item Chat generation wants \texttt{/v1/chat/completions}
    \item Embedding search wants \texttt{/v1/embeddings}
\end{itemize}

They may also need different models, different ports, different throughput assumptions, and
different memory budgets. That is why a dedicated embedding server is usually cleaner than
trying to make one chat model do everything.

%%%%%%%%%%%%%%%%%%%%%%%%%%%%%%%%%%%%%%%%%%%%%%%%%%%%%%%%%%%%%%%%%%%%%%%%%%%
\section{What Gets Indexed}
%%%%%%%%%%%%%%%%%%%%%%%%%%%%%%%%%%%%%%%%%%%%%%%%%%%%%%%%%%%%%%%%%%%%%%%%%%%

In a retrieval workflow, you usually do not embed the entire corpus at query time. Instead:

\begin{enumerate}
    \item Split documents into chunks
    \item Run \texttt{llama-embedding} or an embedding server over those chunks once
    \item Store the vectors in a vector DB or ANN index
    \item At query time, embed only the query
    \item Compare query vector to stored document vectors
\end{enumerate}

This is what makes embedding retrieval scalable.

%%%%%%%%%%%%%%%%%%%%%%%%%%%%%%%%%%%%%%%%%%%%%%%%%%%%%%%%%%%%%%%%%%%%%%%%%%%
\section{Why Normalization Matters}
%%%%%%%%%%%%%%%%%%%%%%%%%%%%%%%%%%%%%%%%%%%%%%%%%%%%%%%%%%%%%%%%%%%%%%%%%%%

Vector similarity is often computed with cosine similarity, dot product, or inner product.
Many retrieval systems either explicitly normalize vectors or assume a similarity convention
that effectively depends on normalization behavior.

The practical takeaway: keep the indexing and query pipelines consistent. Do not change vector
post-processing halfway through a deployment. If query vectors and document vectors are
produced differently, retrieval quality can silently degrade.

%%%%%%%%%%%%%%%%%%%%%%%%%%%%%%%%%%%%%%%%%%%%%%%%%%%%%%%%%%%%%%%%%%%%%%%%%%%
\section{Why Embedding Performance Differs from Generation Performance}
%%%%%%%%%%%%%%%%%%%%%%%%%%%%%%%%%%%%%%%%%%%%%%%%%%%%%%%%%%%%%%%%%%%%%%%%%%%

Embedding workloads are more like: encode text, return one vector. Generation workloads are:
repeatedly decode tokens step by step.

So a model that feels fine for embedding may be much too slow for chat generation, and vice
versa. That is another reason embedding-specialized models are valuable.

%%%%%%%%%%%%%%%%%%%%%%%%%%%%%%%%%%%%%%%%%%%%%%%%%%%%%%%%%%%%%%%%%%%%%%%%%%%
\section{Practical llama-server Invocations}
%%%%%%%%%%%%%%%%%%%%%%%%%%%%%%%%%%%%%%%%%%%%%%%%%%%%%%%%%%%%%%%%%%%%%%%%%%%

\textbf{Start an embedding server with Qwen3-Embedding-4B-GGUF:}

\begin{Verbatim}
llama-server \
  --model artifacts/Qwen3-Embedding-4B-Q4_K_M.gguf \
  --embedding \
  --pooling last \
  --port 8081 \
  --ctx-size 4096
\end{Verbatim}

\textbf{Verify the endpoint is live:}

\begin{Verbatim}
curl -s http://localhost:8081/v1/embeddings \
  -H "Content-Type: application/json" \
  -d '{"input": "test query", "model": "qwen3-embedding"}' \
  | python3 -m json.tool | head -20
\end{Verbatim}

\textbf{Generate a single embedding from the CLI:}

\begin{Verbatim}
llama-embedding \
  --model artifacts/Qwen3-Embedding-4B-Q4_K_M.gguf \
  --pooling last \
  --prompt "Instruct: Retrieve a relevant passage.
Query: how does KV cache work in llama.cpp?"
\end{Verbatim}

\textbf{Run-server.py integration (via \texttt{MODEL\_SLUG}):}

\begin{Verbatim}
MODEL_SLUG=qwen3-embedding-4b ./run-server.py
\end{Verbatim}

%%%%%%%%%%%%%%%%%%%%%%%%%%%%%%%%%%%%%%%%%%%%%%%%%%%%%%%%%%%%%%%%%%%%%%%%%%%
\section{Practical Deployment Pattern for This Repo}
%%%%%%%%%%%%%%%%%%%%%%%%%%%%%%%%%%%%%%%%%%%%%%%%%%%%%%%%%%%%%%%%%%%%%%%%%%%

If this repo adopts \texttt{Qwen3-Embedding-4B-GGUF}, the most likely operational pattern is:

\begin{enumerate}
    \item Use \texttt{llama-embedding} to validate quality and inspect behavior locally
    \item Serve the model with \texttt{llama-server -{}-embedding -{}-pooling last}
    \item Point Open WebUI or another RAG client at that dedicated embeddings endpoint
    \item Later add a Qwen3-Reranker if top-of-list quality becomes the bottleneck
\end{enumerate}

The embedding model and generation model run as separate \texttt{llama-server} instances,
usually on different ports. This avoids context-switch overhead and allows independent tuning
of each model's memory budget.

\begin{tipbox}[Why the embedding model should come first]
If you only add one model first, the embedding model is the right first step because it gives
you the core \texttt{/v1/embeddings} capability, is the prerequisite for scalable corpus
search, and Open WebUI and RAG systems can immediately benefit from it. Generation quality
can be improved later by upgrading the chat model; retrieval quality requires the embedding
model to exist at all.
\end{tipbox}
